\section{Coupled Degenerate System}
Consider a degenerate scaled darcy-stokes coupled boundary-value problem:

Find velocity-pressure potential pair $(\check{\mathbf{v}}_r, \check{q}_l)$ and 
$(\check{\mathbf{v}}_s, \check{q})$ such that 
\begin{alignat}{2}
\label{eq:dimensionlessdarcy_flow_sc_bv}
		  &\check{\mathbf{v}}_r
		  + \phi^{1+\Theta}\nabla(\phi^{-1/2}\check{q}_l)
        &&= 0,\\
\label{eq:dimensionlessconsMass_flow_sc_bv}
		  &\phi^{-1/2}\nabla \cdot(\phi^{1+\Theta}\check{\mathbf{v}}_r)
        + \frac{1}{1-\phi}(\check{q}_l - \phi^{1/2}\check{q}) &&=0,\\
\label{eq:dimensionlessstokes_flow_sc_bv}
		  &\nabla \check{q} - \nabla \cdot \check{\tau}(\check{\mathbf{v}}_s) &&= -(1-\phi)\mathbf{n},\\
\label{eq:dimensionlesscompaction_flow_sc_bv}
		  &\nabla \cdot \check{\mathbf{v}}_s
        - \frac{\phi^{1/2}}{1-\phi}(\check{q}_l - \phi^{1/2}\check{q}) &&=0,
\end{alignat}
and
\begin{equation}\label{eq:bv}
		  \begin{aligned}
					 \check{\mathbf{v}}_r &= \mathbf{g}_{r} \hspace{0.5cm} \text{on} \hspace{0.2cm} \Gamma_{u,D}, \\
					 \check{\mathbf{v}}_s &= \mathbf{g}_{s} \hspace{0.5cm} \text{on} \hspace{0.2cm} \Gamma_{u,S}, \\
					 p &= p_0 \hspace{0.5cm} \text{on} \hspace{0.2cm} \Gamma_{i,D}, \\
					 \pmb{\tau}\cdot \pmb{\nu} - p &= \pmb{t}_0 \hspace{0.56cm} \text{on} \hspace{0.2cm} \Gamma_{i,S},
		  \end{aligned}
\end{equation}
where $\Gamma_{u,D}$ and $\Gamma_{u,S}$ denote the part of boundary we assign essential boundary conditions for 
Darcy and Stokes problem respectively, 
and $\Gamma_{i,D}$ and $\Gamma_{i,S}$ represent the part of boundary we assign natural boundary conditions 
respectively. We have $\Gamma_{u,D} \cup \Gamma_{i,D} = \partial \Omega$ and $\Gamma_{u,S} \cup \Gamma_{i,S} = \partial \Omega$.

We follow the procedure by~\cite{DegenerateDarcy} to discretize this boundary value problem.
To begin weith, a new Hilbert space is defined for this 
specific scaled formulation, where
the scaled velocity $\check{\mathbf{v}}_r$ should lie in, so
\begin{equation}\label{eq:scaledHspace}
		  \mathbb{V}^{\phi}_d = H_{\phi}(\text{div}; \Omega) = \Big\{ \mathbf{v} 
		  \in (L^2(\Omega))^2: \phi^{-1/2}\nabla \cdot (\phi^{1+\Phi} \mathbf{v})
		  \in L^2(\Omega) \Big\},
\end{equation}
which equipped with the inner product as
\begin{equation}\label{eq:scaledinnerp}
		  (\mathbf{u}_{\phi}, \mathbf{v}_{\phi})_{\mathbb{V}^{\phi}_r} = 
		  (\mathbf{u}_{\phi}, \mathbf{v}_{\phi}) + 
		  (\phi^{-1/2} \nabla \cdot (\phi^{1+\Theta}\mathbf{u}_\phi), \phi^{-1/2} \nabla \cdot (\phi^{1+\Theta}\mathbf{v}_\phi)).
\end{equation}
The normal trace on $\partial \Omega$ is well-defined as well as
\begin{equation}\label{eq:tracescaled}
		  \norm{\phi^{1/2+\Theta}\mathbf{v}_\phi\cdot \nu}_{-1/2, \partial \Omega} < C_{\Omega} \Big\{ \norm{\mathbf{v}_\phi} + \norm{\phi^{-1/2}\nabla \cdot (\phi^{1+\Theta}\mathbf{v}_\phi)} \Big\}.
\end{equation}

\subsection{The weak formulation}
We define the following function spaces
\begin{equation}\label{eq:spaces}
		  \begin{aligned}
					 \mathbb{V}^\phi_D &= \Big\{\mathbf{v}\in H_{\phi}(\text{div};\Omega): \phi^{1/2+\Theta} \mathbf{v}\cdot \pmb{\nu} = 0 \hspace{0.15cm} \text{on}\hspace{0.15cm} \Gamma_{u,D}  \Big\},\\
					 \mathbb{W}_D      &= L^2(\Omega),\\
					 \mathbb{V}_S &= \Big\{\mathbf{v}\in (H^1(\Omega))^2 : \mathbf{v} = \mathbf{0} \hspace{0.15cm} \text{on}\hspace{0.15cm}  \Gamma_{u,S} \Big\},\\
		  \mathbb{W}_S &= L^2(\Omega),
		  \end{aligned}
\end{equation}
where each space is equipped with their natural norm.
To empose essential boundary conditoins, we assume that a function $\mathbf{g}_s = \check{\mathbf{v}}_{s,0} \hspace{0.2cm} \text{on} \hspace{0.2cm} \Gamma_{u,S}$ 
and $\mathbf{g}_s = \mathbf{0} \hspace{0.2cm} \text{on} \hspace{0.2cm} \partial \Omega / \Gamma_{u,S}$.
Then $\mathbf{g}_s \in (H^{1/2}(\partial \Omega))^2$ and
extended it continuously from the boundary to the domain, so that
the extension $\mathbf{g}_{s} \in (H^1(\Gamma_{u,S}))^2$ and $\norm{\mathbf{g}_{s,0}}_1 \leq C \norm{\mathbf{v}_{s,0}}_{1/2,\Gamma_{u,S}}$. 
Similarly, we can define a bounded extension of $\check{\mathbf{v}}_r \in \mathbb{V}_D^\phi$ such that                                                                                                                                                 

Scaled nondimensionalized weak form.
Find $\check{\mathbf{v}}_r \in \mathbb{V}^{\phi}_D + \tilde{\mathbf{g}}_{r}$,
$\check{q}_l \in \mathbb{W}_D$, $\check{\mathbf{v}}_s \in \mathbb{V}_S + \tilde{\mathbf{g}}_S$, and $\check{q}\in \mathbb{W}_S$ such that
\begin{alignat}{2}
		  &(\check{\pmb{\tau}}(\bold{v}_{s}),\nabla \pmb{\psi}_s) - (\check{q}, \nabla \cdot \pmb{\psi}_s) = -((1-\phi)\bold{n}, \pmb{\psi}_s) &&\forall\pmb{\psi}_s \in \mathbb{V}_S,\\
					 &(\nabla \cdot \check{\bold{v}}_{s}, \omega)  + \Big(\frac{\phi^{1/2}}{1-\phi} (\check{q}_l - \phi^{1/2}\check{q}), \omega\Big) = 0 &&\forall \omega \in \mathbb{W}_S,\\
					 &(\bold{\check{v}}_{r},\pmb{\psi}_r)  - (\check{q}_{l},\phi^{-1/2} \nabla \cdot (\phi^{1+\Theta}\check{\bold{v}}_r)) = 0, &&\forall \pmb{\psi}_r \in \mathbb{V}_D^\phi,\\
					 \nonumber
					 &(\phi^{-1/2} \nabla \cdot (\phi^{1+\Theta}\check{\bold{v}}_r),\omega_f) + \Big(\frac{1}{1-\phi}(\check{q}_{l} - \phi^{1/2}\check{q}),\omega_f \Big) &&= 0 \\
		  &					 &&\forall \omega_f \in 
		  \mathbb{W}_D.
\end{alignat}
The Existence and uniquense of the solution has been proved by~\cite{DegenerateDarcy}. We will not repeat the proof here.

\subsection{MFEM discretization}
We discretize the formed degenerate system with previously 
defined $H(\text{div})$ and $H^1$ conforming space.
We write the function spaces out as
\begin{equation}\label{eq:femspaces}
		  \begin{aligned}
					 \mathbb{V}_1^0 &= (\mathbb{P}_1)^2 \oplus 
		  \text{curl}\mathbb{S}_2^{\mathcal{DS}}\\
					 \mathbb{W} &= \mathbb{P}_0 \\
					 \mathbb{V}^{\mathcal{BR}}_1 &= \mathcal{DS}_1 \oplus \text{span}\{\mathbf{b}_0,\mathbf{b}_1,\mathbf{b}_2,\mathbf{b}_3 \} \\
					 \mathbb{W}   &= \mathbb{P}_0.
		  \end{aligned}
\end{equation}
We can impose essential boundary conditions

Scaled mixed finite element method. Find $\check{\mathbf{v}}_{r,h} \in \mathbb{V}_1^0 + \hat{\mathbf{g}}_r$, $\check{q})_{r,h} \in \mathbb{W}_D$, 
$\mathbf{v}_{s,h} \in \mathbf{V}_S + \hat{\mathbf{g}}_s$ and $\check{q}_h \in \mathbb{W}_S$ such that
\begin{alignat}{2}
		  \label{eq:weak_stokes}
		  &(\check{\pmb{\tau}}(\bold{v}_{s,h}),\nabla \pmb{\psi}_s) - (\check{q}_h, \nabla \cdot \pmb{\psi}_s) = -((1-\phi)\bold{n}, \pmb{\psi}_s) &&\forall\pmb{\psi}_s \in \mathbb{V}_1^{\mathcal{BR}},\\
		  \label{eq:weak_compaction}
		  &(\nabla \cdot \check{\bold{v}}_{s,h}, \omega)  - \Big(\frac{\phi^{1/2}}{1-\phi} (\check{q}_{l,h} - \phi^{1/2}\check{q}_h), \omega\Big) = 0 &&\forall \omega \in \mathbb{W}_0^{\mathcal{BR}},\\
					\label{eq:weak_darcy}
					 &(\bold{\check{v}}_{r,h},\pmb{\psi}_r)  - (\check{q}_{l,h},\phi^{-1/2} \nabla \cdot (\phi^{1+\Theta}\pmb{\psi}_{r,h})) = 0, &&\forall \pmb{\psi}_r \in \mathbb{V}_1^0,\\
					 \nonumber
					 &(\phi^{-1/2} \nabla \cdot (\phi^{1+\Theta}\check{\bold{v}}_{r,h}),\omega_f) + \Big(\frac{1}{1-\phi}(\check{q}_{l,h} - \phi^{1/2}\check{q}_h),\omega_f \Big) &&= 0 \\
					 \label{eq:weak_continuous}
		  &					 &&\forall \omega_f \in 
		  \mathbb{W}_{0}.
\end{alignat}

\subsection{Local Mass Conservation}
Before we continue the discussion, we need to make sure the formulation is 
well-defined by not dividing by zero.
The term $\phi^{-1/2}$ in the symmetric equations~\eqref{eq:weak_darcy} 
and~\eqref{eq:weak_continuous} causes trouble 
when there is no melting, i.e., $\phi = 0$.
In order to deal with it, we first define a element-averaged $\overline{\phi}_E$ 
as
\begin{equation}\label{eq:averaged_phi}
		  \overline{\phi}_E = \frac{1}{\lvert E \rvert} \int_E \phi(\mathbf{x}) d\mathbf{x},
\end{equation}
where $\lvert E \rvert$ is the area of element $E$.
Let
\begin{equation}\label{eq:newvariable_phi}
		  \phi_E = \left\{
					 \begin{aligned}
								&\overline{\phi}_E  \hspace{0.2cm} \text{if}\hspace{0.2cm} \overline{\phi}_E \neq 0,\\
								&0 \hspace{0.5cm} \text{if} \hspace{0.2cm} \overline{\phi}_E = 0.
					 \end{aligned}
					 \right.
\end{equation}
We now replace the $\phi^{-1/2}$ with $\phi_E^{-1/2}$. Then we expand the integral
with divergence theorem as
\begin{equation}\label{eq:divtheorem}
		  \begin{aligned}
					 \int_E \phi_E^{-1/2} \nabla \cdot (\phi^{1+\Theta} \pmb{\psi}_{r,h})\omega_l &=\int_{\partial E} \phi_E^{-1/2} \phi^{1+\Theta} \pmb{\psi}_{r,h}\cdot \pmb{\nu} \omega_f ds \\
					 &-\int_{E}\phi_E^{-1/2}\phi^{1+\Theta}\pmb{\psi}_{r,h} \cdot \nabla \omega_l d\mathbf{x}.
		  \end{aligned}
\end{equation}
The second term vanishes when $\omega_l$ is constant on element $E$.
We define an auxiliary auxiliary space $\omega_E$ begin 1 on $E$ and 0 elsewhere.
We sum equations~\eqref{eq:weak_compaction} and~\eqref{eq:weak_continuous} 
together 
assuming that $\omega = \omega_E \in \mathbb{W}_{S}$ and $\omega_l = \phi_E^{-1/2} \omega_E \in \mathbb{W}_{D}$. We can achieve local mass conservation of the
phase averaged mixture if we
replace the term $\phi^{1/2}$ in equation~\eqref{eq:weak_compaction} with
$\phi_E^{1/2}$ as
\begin{equation}
		  \int_{E}\nabla \cdot (\phi^{1+\Theta}\check{\mathbf{v}}_r + \check{\mathbf{v}}_s) d\mathbf{x} = 0.
\end{equation}
We now propose our local mass conserved version as
\begin{alignat}{2}
		  \label{eq:weak_stokes_mass}
		  &(\check{\pmb{\tau}}(\bold{v}_{s,h}),\nabla \pmb{\psi}_s) - (\check{q}_h, \nabla \cdot \pmb{\psi}_s) = -((1-\phi)\bold{n}, \pmb{\psi}_s) &&\forall\pmb{\psi}_s \in \mathbb{V}_1^{\mathcal{BR}},\\
		  \label{eq:weak_compaction_mass}
		  &(\nabla \cdot \check{\bold{v}}_{s,h}, \omega)  - \Big(\frac{\phi_E^{1/2}}{1-\phi} (\check{q}_{l,h} - \phi^{1/2}\check{q}_h), \omega\Big) = 0 &&\forall \omega \in \mathbb{W}_0^{\mathcal{BR}},\\
					\label{eq:weak_darcy_mass}
					 &(\bold{\check{v}}_{r,h},\pmb{\psi}_r)  - (\check{q}_{l,h},\phi_E^{-1/2} \nabla \cdot (\phi^{1+\Theta}\pmb{\psi}_{r,h})) = 0, &&\forall \pmb{\psi}_r \in \mathbb{V}_1^0,\\
					 \nonumber
					 &(\phi_E^{-1/2} \nabla \cdot (\phi^{1+\Theta}\check{\bold{v}}_{r,h}),\omega_f) + \Big(\frac{1}{1-\phi}(\check{q}_{l,h} - \phi^{1/2}\check{q}_h),\omega_f \Big) &&= 0 \\
					 \label{eq:weak_continuous_mass}
		  &					 &&\forall \omega_f \in 
		  \mathbb{W}_{0}.
\end{alignat}

\subsection{Convergence Rate}
We derive a bound for the error by taking the difference
of 


\subsection{Test Run}
We test our MFEM setup with two problems.
First, we test with a manufactured solution. Second, we test a
more complex yet a more realistic problem.

Problem setup 1:

We solve a manufactured system of equations of the rectangular domain $-1< x < 1$, 
\begin{equation}\label{eq:test1}
		  \mathbf{v}_s = 
\end{equation}

Problem setup 2:

We solve the full system of equations on the rectangular domain $-160\text{km} < x < 160\text{km}$,
$0, x, 160\text{km}$. The boundary conditions are define by imposing  
\begin{equation}\label{eq:test2}
		  \mathbf{v}_s = 
\end{equation}
